% !TEX root = ../ADEmain.tex
\Section{Manifold learning}

Consider a regression \( Y_{i} = \fs(\Xv_{i}) + \eps_{i} \) with a high-dimensional \( \Xv_{i} \in \R^{\dimd} \),
\( n = 1,\ldots,n \).
The manifold hypothesis assumes that the function \( \fs(\cdot) \) has a nearly multi-index structure in a local vicinity 
of each point \( \xv \in \R^{\dimd} \).
This means a local linear approximation
\begin{EQA}[c]
	\E Y_{i}
	\approx
	\fcn_{j} + \fclv_{j}^{\T} \PEDR_{j} (\Xv_{i} - \xv_{j})
\end{EQA}
for all points \( \Xv_{i} \) in a vicinity of \( \xv_{j} \).
Here not only the intercept \( \fcn_{j} \in \R \) and 
the slope coefficients \( \fclv_{j} = (\fcl_{1},\ldots,\fcl_{\dimm})^{\T} \in \R^{\dimm} \)
but also the mapping \( \PEDR_{j} = \PEDR(\xv) \in \PROJ_{\dimd,\dimm} \)
depend on \( \xv_{j} \).
%
The proposed approach extends the multi-index procedure from Section~\ref{Smultind} to this situation 
by considering a family \( (\PEDR_{j}) \) in place of one mapping \( \PEDR \) with a strong penalty on their variability.
This approach does not require any manifold description.
We also do not project the data on the manifold.

%%%%%%%%%%%%%%%%%%%% section %%%%%%%%%%%%%%%%%%%%
\Subsection{Multiscale modeling}
The manifold hypothesis means that for any fixed \( \xvd \) and any \( \xv \) close to \( \xvd \), 
one can use the mapping \( \PEDR(\xvd) \) in place of \( \PEDR(\xv) \). 
%
This discussion suggests introducing two scales.
One of them is for local linear approximation of the function \( \fs(\cdot) \) in the form
\begin{EQA}[c]
	\fs(\xv) - \fs(\xvd)
	\approx
	\fclv^{\T} \, \PEDR (\xv - \xvd) 
\end{EQA} 
with the coefficients \( \fclv = (\fcl_{m}) \) and mappings \( \PEDR \in \PROJ_{\dimd,\dimm} \) depending
on the central point \( \xvd \).
Another one for local manifold approximation:  
\(
	\PEDR(\xv)^{\T} \PEDR(\xv)
	\approx
	\PEDR(\xvd)^{\T} \PEDR(\xvd)
%	\, .
\).
%A local tensor 
%\( \TEDR_{j} \) and bandwidth \( \hiso_{0} \) for the function \( \fs \),
%and tensor \( \TEDRm_{j} \) and bandwidth \( \hisom \) for the manifold.
At the initialization step of the procedure, we consider an isotropic vicinity
\( \{ \| \xv - \xvd \| \leq \hiso_{0} \} \) for approximation of \( \fs \)
and \( \{ \| \xv - \xvd \| \leq \hisom_{0} \} \) for continuity of \( \PEDR(\xv) \).
The approach assumes that the manifold is ``smoother'' than the function \( \fs \) and hence,
\( \hisom_{0} \gg \hiso_{0} \).

%, where \( \hiso_{0} \) is used for the weights \( \weight_{ij} \).
One can go one step further and assume that the manifold belongs to a linear subspace \( \EDR \) of a larger dimension \( \dimms \).
An example is a smooth curve in a two- or three-dimensional subspace.
Under such an assumption, one can first apply the multi-index procedure to recover this larger subspace \( \EDR \)
and then use this information for a warm start.

%%%%%%%%%%%%%%%%%%%% section %%%%%%%%%%%%%%%%%%%%
\Subsection{Objective function and manifold penalty}
The approach extends the procedure for the multi-index case from Section~\ref{SEDRmiproc}.
The basic idea behind the procedure from \eqref{ju7fudeju74334fhfd} is that the EDR space 
is the same for all \( \xv_{j} \).
This structural assumption does not hold in the manifold case.
However, it is nearly fulfilled if we restrict the analysis to a vicinity of each \( \xv_{j} \).
This suggests using two hierarchical localizing schemes described by 
localizing tensors \( \TEDR \) and \( \TEDRm \).

Suppose a collection of centers \( ( \xv_{j} )_{j \in \JJJ} \), and 
the sets of directions \( (\Sph_{j})_{j \in \JJJ} \) have been fixed.
% 
For each \( j \in \JJJ \), let \( \neff_{j} \),
\( \Imav_{j} = \bigl( \Ima_{j,\sph} \bigr)_{\sph \in \Sph_{j}} \),
%\( \Sv_{j} = \bigl( S_{j,\sph} \bigr)_{\sph \in \Sph_{j}} \),
\( \UV_{j}^{\T} = (\Uv_{j,\sph})_{\sph \in \Sph_{j}} \),
be computed as in \eqref{uyhvvy883hierwfewAVE} and \eqref{juu8hjfioee433erhAVE} with 
\( \weight_{ij} = \Kern(\| \TEDR_{0} (\Xv_{i} - \xv_{j}) \|^{2}) \)
for \( \TEDR_{0} = \hiso_{0}^{-1} \Id_{\dimd} \) and a starting bandwidth \( \hiso_{0} \).
%
Now, for any \( \xv_{\jl} \in (\xv_{j}) \), consider a localized version of
\eqref{ju7fudeju74334fhfd}:
\begin{EQA}[rcl]
%	\prmtv^{(0)} 
	&& \nquad
	\bigl\{ \PEDR_{\jl} \, , (\fclv_{\jjl})_{j \in \JJJ} \bigr\}
	= 
	\arginf_{\PEDR \, , (\fclv_{j})_{j \in \JJJ}} 
	\sum_{j \in \JJJ} 
	\bigl\| \Imav_{j} - \UV_{j} \PEDR^{\T} \fclv_{j} \bigr\|^{2} \neff_{j} \,
	\weightM_{\jjl}
%	\Kern\bigl( \| \TEDRm(\xv_{j} - \xv_{\jl}) \|^{2} \bigr)
%	+ \lambda_{\PEDR} \pent(\BEDR)
	\, ,
	\qquad \quad
\label{kjhviir8786tytcdeefm}
\end{EQA}
where \( \weightM_{\jjl} = \Kern(\| \TEDRm_{0} (\xv_{j} - \xv_{\jl}) \|^{2}) \) for
\( \TEDRm_{0} = \hisom_{0}^{-1} \Id_{\dimd} \) and \( \hisom_{0} \gg \hiso_{0} \).
The only difference between this problem and \eqref{ju7fudeju74334fhfd} is in using a multiplier (localizing weight)
\( \weightM_{\jjl} \)
%\( \Kern\bigl( \| \TEDRm(\xv_{j} - \xv_{\jl}) \|^{2} \bigr) \)
for each entry \( \bigl\| \Imav_{j} - \UV_{j} \PEDR^{\T} \fclv_{j} \bigr\|^{2} \).
Such weighting restricts the multi-index approximation to a vicinity of \( \xv_{\jl} \) given by 
such weights \( (\weightM_{\jjl})_{j \in \JJJ} \). 
All local problems \eqref{kjhviir8786tytcdeefm} can be merged into one:
%leading to minimization of
\begin{EQA}[c]
	\sum_{\jl \in \JJJ} \sum_{j \in \JJJ}
	\bigl\| \Imav_{j} - \UV_{j} \PEDR_{\jl}^{\T} \fclv_{\jjl} \bigr\|^{2} \neff_{j} \, \weightM_{\jjl}
	\to 
	\inf_{\{ \PEDR_{\jl} \, , (\fclv_{\jjl})_{j \in \JJJ} \}_{\jl \in \JJJ}}
	\, .
	\qquad \quad
\label{kjhviir8786tytcdeefde}
\end{EQA}
%over \( \prmtv = \{ \PEDR_{\jl} \, , (\fclv_{\jjl})_{j \in \JJJ} \}_{\jl \in \JJJ} \).
%
The choice of the family \( (\PEDR_{\jl}) \) can be stabilized and synchronized by a manifold penalty
\begin{EQA}[c]
	\pent_{\subM}(\PEDR)
	\eqdef
	\lambda \sum_{\jl \in \JJJ} \sum_{j \in \JJJ} \epsilon(\PEDR_{j} , \PEDR_{\jl})
%	\| \PEDR_{j} - \PEDR_{\jl} \|_{\Fr}^{2} 
	\; \weightM_{\jjl}
%	\Kern\bigl( \| \TEDRm_{j} (\xv_{j} - \xv_{\jl}) \|^{2} \bigr)
	\, .
\end{EQA}
The value \( \epsilon(\PEDR_{j} , \PEDR_{\jl}) \) measures discrepancy between \( \PEDR_{j} \) and \( \PEDR_{\jl} \).
The Frobenius or operator norms of the difference \( \PEDR_{j} - \PEDR_{\jl} \) cannot be used here because
each \( \PEDR_{j} \) is not uniquely defined; one can swap the sign of each row without changing the manifold projection. 
%The difference between projectors \( \PEDR_{j} \PEDR_{j}^{\T} - \PEDR_{\jl} \PEDR_{\jl}^{\T} \) is a reasonable alternative.
The basic idea behind the manifold penalty is that 
for close points \( \xv_{j} \) and  \( \xv_{\jl} \), the corresponding projectors  
\( \PEDR_{j}^{\T} \PEDR_{j} \) and \( \PEDR_{\jl}^{\T} \PEDR_{\jl} \) are also close to each other.  
This suggests
\begin{EQA}[c]
	\epsilon(\PEDR_{j} , \PEDR_{\jl})
	\eqdef
	\bigl\| \PEDR_{j} (\Id_{\dimd} - \PEDR_{\jl}^{\T} \PEDR_{\jl}) \bigr\|_{\Fr}^{2}
	\, .
\end{EQA}
This leads to minimization of %\( \LL_{\subM}(\prmtv) \) with
\begin{EQA}[c]
%	&& \nquad
%	(\BEDR_{\jl} \, , (\fclv_{\jjl})_{j \in \JJJ} )_{\jl \in \JJJ}
%	\\
%	&=& 
%	\arginf_{(\BEDR_{j} \, , \fclv_{\jjl})_{j \in \JJJ}} 
	\LL_{\subM}(\prmtv)
	\eqdef
	\sum_{\jl \in \JJJ} \sum_{j \in \JJJ}
	\bigl\| \Imav_{j} - \UV_{j} \PEDR_{\jl}^{\T} \fclv_{\jjl} \bigr\|^{2} \neff_{j} \, \weightM_{\jjl}
	+ \lambda \sum_{\jl \in \JJJ} \sum_{j \in \JJJ} \epsilon( \PEDR_{j} , \PEDR_{\jl}) \; \weightM_{\jjl}
	\, .
\end{EQA}


%%%%%%%%%%%%%%%%%%%% section %%%%%%%%%%%%%%%%%%%%
\Subsection{Basic (one-step) procedure}
\label{SmaniEDR}
The global problem of minimizing the objective function \( \LL_{\subM}(\prmtv) \)
can be solved by alternating.
We assume that  
a preliminary guess \( \PEDR_{j,\init} \) is available for each point from the family \( (\xv_{j})_{j \in \JJJ} \).
Fix \( \jl \in \JJJ \) and consider a partial optimization problem:
\begin{EQA}[c]
	\sum_{j \in \JJJ}
	\bigl\| \Imav_{j} - \UV_{j} \PEDR_{\jl}^{\T} \fclv_{\jjl} \bigr\|^{2} \neff_{j} \, \weightM_{\jjl}
	+ \lambda \sum_{j \in \JJJ} \epsilon( \PEDR_{\jl} , \PEDR_{j,\init}) \; \weightM_{\jjl}
	\to \inf_{\PEDR_{\jl} \, , (\fclv_{\jjl})_{j \in \JJJ}}
	\, .
	\qquad
\label{iuhviiioijcwtyuvb}
\end{EQA}
It can be solved by the procedure from Section~\ref{SaltMI}
with obvious changes.
Given \( \PEDR_{\jl,\init} \), each coefficient \( \fclv_{\jjl} \) for \( j \in \JJJ \) can be obtained by soving 
the quadratic problem
\begin{EQA}[c]
	\bigl\| \Imav_{j} - \UV_{j} \PEDR_{\jl,\init}^{\T} \fclv_{\jjl} \bigr\|^{2}
	\to \inf_{\fclv_{\jjl}}
	\, .
	\qquad
\label{kiiilfduy543rbggioe46}
\end{EQA}
Minimization of {\color{blue} \eqref{iuhviiioijcwtyuvb}} w.r.t. \( \PEDR_{\jl} \) is a bit more involved,
a relaxation of the constraint \( \PEDR_{\jl} \in \PROJ_{\dimd,\dimm} \) is helpful.
Define
\begin{EQA}[c]
	\IEDR_{\jl,\init}
	\eqdef
	\frac{1}{\WeightM_{\jl}} \sum_{j \in \JJJ} \PEDR_{j,\init}^{\T} \, \PEDR_{j,\init} \; \weightM_{\jjl}
	\, ,
	\qquad
	\WeightM_{\jl}
	\eqdef
	\sum_{j \in \JJJ} \weightM_{\jjl}	
	\, .
\label{jfu88uuujecvyyyyyue}
\end{EQA}
Note that \( \IEDR_{\jl,\init} \) is not a projector, because the set of projectors is not convex.
However, it is a subprojector in the sense that \( 0 \leq \IEDR_{\jl,\init} \leq \Id_{\dimd} \).



\begin{lemma}
\label{Lmanijjl}
With \( (\fclv_{\jjl})_{j \in \JJJ} \) fixed and \( \IEDR_{\jl,\init} \) from \eqref{jfu88uuujecvyyyyyue}, it holds
%the solution to \eqref{iuhviiioijcwtyuvb} coincides with the solution to
\begin{EQA}[rcl]
	&& \nquad
	\arginf_{\BEDR_{\jl} \in \R^{\dimm \times \dimd}} \biggl( 
		\sum_{j \in \JJJ} \bigl\| \Imav_{j} - \UV_{j} \BEDR_{\jl}^{\T} \fclv_{\jjl} \bigr\|^{2} \neff_{j} \, \weightM_{\jjl}
		+ \lambda \sum_{j \in \JJJ} 
			\bigl\| \BEDR_{\jl} (\Id_{\dimd} - \PEDR_{j,\init}^{\T} \PEDR_{j,\init}) \bigr\|_{\Fr}^{2} \; \weightM_{\jjl} 
	\biggr) 
	\\
	&=&
	\arginf_{\BEDR_{\jl} \in \R^{\dimm \times \dimd}} 
	\biggl( \sum_{j \in \JJJ}
		\bigl\| \Imav_{j} - \UV_{j} \BEDR_{\jl}^{\T} \fclv_{\jjl} \bigr\|^{2} \neff_{j} \, \weightM_{\jjl}
		+ \lambda \WeightM_{\jl} \tr \bigl\{ \BEDR_{\jl}^{\T} \BEDR_{\jl} (\Id_{\dimd} - \IEDR_{\jl,\init} ) \bigr\}
	\biggr)
	\, .
	\qquad
\label{iuhviiioijcwtyuvbr}
\end{EQA}
\end{lemma}

\begin{proof}
{\color{blue} please check using the following lemma}
\end{proof}

\begin{lemma}
\label{LFrobML}
For any positive weights \( \weightM_{j} \) and projectors \( \PEDR_{j} \in \PROJ_{\dimd,\dimm} \) for \( j \in \JJJ \), it holds
\begin{EQA}[c]
	\sum_{j \in \JJJ} \| \BEDR (\Id_{\dimd} - \PEDR_{j}^{\T} \PEDR_{j}) \|_{\Fr}^{2} \, \weightM_{j}
	=
	\WeightM \tr\bigl\{ \BEDR^{\T} \BEDR (\Id_{\dimd} - \IEDR) \bigr\}
	\, ,
\end{EQA}
where
\begin{EQA}[c]
	\WeightM
	\eqdef
	\sum_{j \in \JJJ} \weightM_{j}
	\, ,
	\qquad
	\IEDR
	\eqdef
	\frac{1}{\WeightM} \sum_{j \in \JJJ} \PEDR_{j}^{\T} \PEDR_{j} \, \weightM_{j}
\end{EQA}
\end{lemma}

\begin{proof}
As \( \PEDR_{j}^{\T} \PEDR_{j} \) is a projector in \( \R^{\dimd} \) and\( \PEDR_{j} \, \PEDR_{j}^{\T} = \Id_{\dimm} \), 
 it holds 
\begin{EQA}[rcl]
	&& \nquad
	\sum_{j \in \JJJ} \| \BEDR (\Id_{\dimd} - \PEDR_{j}^{\T} \PEDR_{j}) \|_{\Fr}^{2}	 \, \weightM_{j}
	=
	\sum_{j \in \JJJ}\tr (\BEDR - \BEDR \, \PEDR_{j}^{\T} \PEDR_{j}) (\BEDR - \BEDR \, \PEDR_{j}^{\T} \PEDR_{j})^{\T} \weightM_{j}
	\\
	&=&
	\sum_{j \in \JJJ} \tr\Bigl\{ (\BEDR \, \BEDR^{\T} - \tr(\BEDR^{\T} \BEDR \, \PEDR_{j}^{\T} \PEDR_{j}) \Bigr\} \weightM_{j}
	=
	\WeightM \tr (\BEDR^{\T} \BEDR ) - \WeightM \tr (\BEDR^{\T} \BEDR \, \IEDR)
	\, .
\end{EQA}
This implies the assertion.
\end{proof}


Based on the solution \( \BEDR_{\jl} \) of \eqref{iuhviiioijcwtyuvbr}, compute
\begin{EQA}[rcl]
	\JEDR_{\jl}
	\eqdef
	\sum_{j \in \JJJ} (\BEDR_{\jl}^{\T} \, \fclv_{\jjl}) \, \bigl( \BEDR_{\jl}^{\T} \, \fclv_{\jjl} \bigr)^{\T}
		\neff_{j} \, \weightM_{\jjl}
	&=&
	\BEDR_{\jl}^{\T} \biggl( 
		\sum_{j \in \JJJ} \fclv_{\jjl} \, \fclv_{\jjl}^{\T} \, \neff_{j} \, \weightM_{\jjl} 
	\biggr) \BEDR_{\jl}
	\, .
\label{jhuifdyu8v8i6445hve}
\end{EQA}
As \( \BEDR_{\jl} \) is of rank \( \dimm \), the matrix \( \JEDR_{\jl} \) is 
also a rank \( \dimm \) matrix. 
Compute also SVD 
\begin{EQA}[c]
	\frac{1}{\| \JEDR_{\jl} \|} \JEDR_{\jl}
	=
	\PEDR_{\jl}^{\T} \, \Lambda_{\jl} \, \PEDR_{\jl}
	\, ,
	\qquad
	\PEDR_{\jl} \, \PEDR_{\jl}^{\T}
	= 
	\Id_{\dimm}
	\, ,
	\qquad
	\Lambda_{\jl} = \diag(\lambda_{1} \geq \ldots \geq \lambda_{\dimm})
	\, 
	\qquad
\label{jhfuui887678ucfhie3}
\end{EQA}
with \( \lambda_{1} = 1 \).
%and define \( \PEDR_{\jl,\outp} = \PEDR_{\jl} \) and \( \Lambda_{\jl,\outp} \eqdef \Lambda_{\jl} / \| \Lambda_{\jl} \| \).
Matrices \( \PEDR_{\jl} \, , \Lambda_{\jl} \) describe the structure of the models and will be used 
for the alternating procedure by
replacing \( \PEDR_{\jl,\init} \) with \( \PEDR_{\jl} \).

Now we can describe the \emph{one-step procedure}. 
Assume available:
%
\begin{mylist}
	\item 
	a collection of mappings \( (\PEDR_{\jl,\init})_{\jl \in \JJJ} \) from 
	initialization or the previous step;

	\item 
	a collection \( (\neff_{j}, \Imav_{j} \, , \UV_{j})_{j \in \JJJ} \); 
	see \eqref{uyhvvy883hierwfewAVE} and \eqref{juu8hjfioee433erhAVE};

	\item 
%	manifold tensors \( \TEDRm_{\jl} \) 
%	and manifold weights \( \weightM_{\jjl} = \Kern(\| \TEDRm_{\jl} (\xv_{j} - \xv_{\jl}) \|^{2}) \); 
 	manifold weights \( \weightM_{\jjl} \).
\end{mylist}
%
\paragraph{One-step procedure.}
For each \( \jl \in \JJJ \):
%
\begin{mylist}
	\item
	Fix \( (\fclv_{\jjl})_{j \in \JJJ} \) by solving \eqref{kiiilfduy543rbggioe46}.
%
	\item
	Compute \( \IEDR_{\jl,\init} \) by \eqref{jfu88uuujecvyyyyyue}.
%
	\item
	Solve
\begin{EQA}[c]
	\BEDR_{\jl}
	=
	\arginf_{\BEDR \in \R^{\dimm \times \dimd}} \biggl( \sum_{j \in \JJJ}
		\bigl\| \Imav_{j} - \UV_{j} \BEDR^{\T} \fclv_{\jjl} \bigr\|^{2} \neff_{j} \, \weightM_{\jjl}
		+ \lambda_{\subM} \tr \bigl\{ \BEDR^{\T} \BEDR (\Id_{\dimd} - \IEDR_{\jl,\init}) \bigr\} 
	\biggr)
	\, .
	\qquad
\end{EQA}
%
	\item
	Compute \( \PEDR_{\jl} \) and \( \Lambda_{\jl} \) by  \eqref{jhuifdyu8v8i6445hve}, \eqref{jhfuui887678ucfhie3}.
	Update \( \PEDR_{\jl,\init} = \PEDR_{\jl} \).

\end{mylist}

%%%%%%%%%%%%%%%%%%%% section %%%%%%%%%%%%%%%%%%%%
\Subsection{Initialization by local linear estimation}
%The manifold learning procedure can be initialized using local linear gradient estimation.
Suppose a collection \( (\xv_{j})_{j \in \JJJ} \) to be fixed.
For each \( j \in \JJJ \), we can estimate the gradient \( \nabla \gs(\xv_{j}) \) 
by a local linear fit as suggested in Section~\ref{Siniloclin}.
With a proper choice of the parameter \( \nmin \), \( \nlin \), and \( \nman \):
fix \( \hiso_{0} \), \( \hiso_{\lin} \), and \( \hisom_{0} \) as the smallest values ensuring 
\begin{EQ}[rcl]
	\frac{1}{\nJJJ}\sum_{j \in \JJJ} 
	\sumi \Kern\Bigl( \mfrac{\| \Xv_{i} - \xv_{j} \|^{2}}{\hiso_{0}^{2}} \Bigr)
	& \geq &
	\nmin
	\, ,
	\quad
	\frac{1}{\nJJJ} \sum_{j \in \JJJ} 
	\sumi \Kern\Bigl( \mfrac{\| \Xv_{i} - \xv_{j} \|^{2}}{\hiso_{\lin}^{2}} \Bigr)
	\geq 
	\nlin
	\, ,
	\\
	\frac{1}{\nJJJ} \sum_{\jl \in \JJJ} 
	\sum_{j \in \JJJ} \Kern\Bigl( \mfrac{\| \xv_{j} - \xv_{\jl} \|^{2}}{\hisom_{0}^{2}} \Bigr)
	& \geq &
	\nman
	\, .
\label{hhfiecredgvhihiue86tML}
\end{EQ}
For each \( j \in \JJJ \), compute \( \hat{\nabla \gs}(\xv_{j}) \) by \eqref{ihvifvityguwhd8354e}
%\begin{EQA}[rcl]
%	\hat{\nabla \gs}(\xv_{j})
%	&=&
%	\biggl\{ \sumi (\Xv_{i} - \Xbar_{j}) \, (\Xv_{i} - \Xbar_{j})^{\T}
%		\Kern\Bigl( \mfrac{\| \Xv_{i} - \xv_{j} \|^{2}}{\hiso_{\lin}^{2}} \Bigr)
%	\biggr\}^{-1}
%	\\
%	&&
%	\times
%	\sumi Y_{i} \, (\Xv_{i} - \Xbar_{j}) \, 
%		\Kern\Bigl( \mfrac{\| \Xv_{i} - \xv_{j} \|^{2}}{\hiso_{\lin}^{2}} \Bigr)
%	\, 
%\label{chgvudimmmkkjgcthe}
%\end{EQA}
with \( \Xbar_{j} \) from \eqref{uyhvvy883hierwfewAVE}.
Further, for each \( \jl \in \JJJ \), compute \( \JEDR_{\jl} \) as
\begin{EQA}[c]
	\JEDR_{\jl}
	\eqdef
	\sum_{j \in \JJJ} \hat{\nabla \gs}(\xv_{j}) \bigl\{ \hat{\nabla \gs}(\xv_{j}) \bigr\}^{\T} \neff_{j} \, 
	\Kern\Bigl( \mfrac{\| \xv_{j} - \xv_{\jl} \|^{2}}{\hisom_{0}^{2}} \Bigr)
	\, ,
\label{jheew3regdw3bvio2ff}
\end{EQA}
and initialize \( \PEDR_{\jl,\init} \) by a PCA procedure using \( \dimm \) principal eigenvectors of \( \JEDR_{\jl} \).  
%Suppose that the points \( \xv_{\jl} \) are ordered with a possibly small distance between 
%any two subsequent points.
%The idea is to start with a non-informative guess for the first point \( \xv_{\jl} \) and solve \eqref{kjhviir8786tytcdeefm} 
%using the procedure from Section~\ref{SaltMI}. 
%The solution will be used as \( \PEDR_{\jl}^{(0)} \).
%Then we take the second point and solve \eqref{iuhviiioijcwtyuvb} using already obtained solutions \( \PEDR_{\jl}^{(0)} \)
%for manifold synchronization.
%The procedure should be continued until \( \PEDR_{\jl}^{(0)} \) is initialized for all \( (\xv_{\jl})_{\jl \in \JJJ} \).


%%%%%%%%%%%%%%%%%%%% section %%%%%%%%%%%%%%%%%%%%
\Subsection{Structure-adaptive manifold learning}
\label{SEDRmanproc}
This section describes the manifold learning procedure starting from initialization.




\begin{table}[h]
\begin{tabular}{c|c|c}
        \hline
        Parameter & Meaning & possible values and default \\
        \hline
        \( \nmin \) & \( \# \) of design points \( \Xv_{i} \) in a local vicinity & \( 7,{10}, \bb{15}, 20 \) \\
        \( \nJJJ \) & size of \( \JJJ \) 				& \( s \, n/\nmin \) for \( s=1,\bm{2},\ldots,\nmin \) \\
		\( \nSph \) & size of each \( \Sph_{j} \)		& \( \bm{\nmin}, 2\nmin, 3\nmin \) \\
%\( a (= 2^{1/\dimm}) \), 
%		\( \lambda & (= \nmin \dimm) \),
		\( \lambda_{\subM} \) & Lagrange multiplier 		& \( \bm{1}, 2, \ldots, \nmin \) \\
		\( \Kern \) & the kernel & Epanechnikov: \( \Kern(t) = (1-t^{2})_{+} \) \\
		\( k_{\max} \) & the number of repetitions in AO & 3, \textbf{5}, 7 \\
		\( a \)	& decrease factor for \( \hiso_{k} \) & \( 2^{1/(\dimm+2)}, {2^{1/(\dimm+1)}}, \bm{2^{1/\dimm}} \) \\
		\( \hiso_{\min} \) & the smallest \( \hiso_{k} \) & \( s \sigmaX/\sqrt{n} \, \) for \( s=1,2,\bb{3} \) \\
		\hline
    \end{tabular}
\caption{Tuning parameters for the manifold learning procedure}    
\label{TdefML}
\end{table}

%\begin{mylist}
%\item
%\( \dimm (=1,2,3) \), 
%\item
%\( \nmin (= 10) \), 
%\item
%\( \nJJJ (= n/2) \), 
%\item
%\( \nman (= \nJJJ/10) \),
%\item
%\( \nSph (= \min\{ \nmin, \dimd \}) \),
%%\( a (= 2^{1/\dimm}) \), 
%\item
%\( \lambda (= \nmin \dimm) \),
%\item
%\( \lambda_{\subM} (=\nmin) \).
%\end{mylist}

\paragraph{Initialization.}
%Make preliminary calculations as follows:
\begin{mylist}

\item
Fix the set of tuning parameters; see Table~\ref{TdefML}.

\item
Generate the set \( (\xv_{j})_{j \in \JJJ} \).

\item
Fix \( \hiso_{\lin} \) by \eqref{hhfiecredgvhihiue86tML}.
%as the smallest values ensuring
%\begin{EQA}[rcl]
%	\frac{1}{\nJJJ} \sum_{j \in \JJJ} 
%	\sumi \Kern\bigl( \hiso_{0}^{-2} \| \Xv_{i} - \xv_{j} \|^{2} \bigr)
%	& \geq &
%	\nmin
%	\, .
%\end{EQA}


\item
For each \( j \in \JJJ \), compute \( \hat{\nabla \gs}(\xv_{j}) \) by \eqref{ihvifvityguwhd8354e}.

\item 
For each \( \jl \in \JJJ \), compute \( \JEDR_{\jl} \) by \eqref{jheew3regdw3bvio2ff}
and define \( \PEDR_{\jl,\init} \) 
using the first \( \dimm \) principal directions of \( \JEDR_{\jl} \).

\item
Syncronize the initialization by 
running the one-step procedure several times (up to \( k_{\max} \)) 
using the output \( \PEDR_{\jl} \)
of the previous step as \( \PEDR_{\jl,\init} \) for all \( \jl \in \JJJ \).

\item
Fix \( \hiso_{0} \) as the smallest value ensuring \eqref{jhgiverkytde4w34etycj}.
Define \( \TEDR_{0} = \hiso_{0}^{-1} \Id_{\dimd} \).

\item Set \( k=0 \).

\end{mylist}
%After initialization, it is useful to carry out the one-step procedure several times (up to \( k_{\max} \))
%to make the first-step results more stable.
%After each run, the output \( (\PEDR_{\jl,\init})_{\jl \in \JJJ} \) should be used as the previous step collection 
%\( (\PEDR_{\jl}^{(0)})_{\jl \in \JJJ} \) for the next run.


%\paragraph{Sequential initialization.}
%For a given ordering on \( \JJJ \),
%compute sequentially the initial mappings \( \PEDR_{\jl}^{(0)} \), one after another, as follows:
%\begin{mylist}
%
%\item 
%Take the first point \( \xv_{\jl} \):
%\begin{mylist}
%
%\item
%Initialize \( \PEDR_{\jl}^{(0)} \) randomly or by the first-step global multi-index procedure, 
%set \( \BEDR_{\jl}^{(0)} = \PEDR_{\jl}^{(0)} \).
%
%\item With \( \BEDR_{\jl}^{(0)} \) fixed, apply the localized multi-index procedure to solve \eqref{kjhviir8786tytcdeefm}:
%\begin{mylist}
%
%	\item
%	Fix \( (\fclv_{\jjl})_{j \in \JJJ} \) by solving \eqref{kiiilfduy543rbggioe46}.
%
%	\item
%	Fix \( \BEDR_{\jl} \) as a solution of
%\begin{EQA}[c]
%	\nquad
%%	\nquad
%%	\BEDR_{\jl}
%%	=
%	\inf_{\BEDR \in \R^{\dimm \times \dimd}} \biggl( \sum_{j \in \JJJ}
%		\bigl\| \Imav_{j}^{(0)} - \UV_{j}^{(0)} \BEDR^{\T} \fclv_{\jjl} \bigr\|^{2} \;
%	\weightM_{\jjl}^{(0)}
%		+ \lambda_{\subM} \| \BEDR  - \BEDR_{\jl}^{(0)} \|_{\Fr}^{2} 
%	\biggr)
%	\, .
%	\qquad
%\label{jjhidkmnvuy5rrtvnv}
%\end{EQA}
%
%	\item
%	Compute \( \PEDR_{\jl} \) and \( \Lambda_{\jl} \) by SVD of \( \JEDR_{\jl} \) as in \eqref{jhuifdyu8v8i6445hve} and
%	\eqref{jhfuui887678ucfhie3}.
%
%	\item
%	Update \( \PEDR_{\jl}^{(0)} = \BEDR_{\jl}^{(0)} = \PEDR_{\jl} \) and \( \Lambda_{\jl}^{(0)} = \Lambda_{\jl} \).
%	
%	\item 
%	Loop \( k_{\max} \) times.
%	
%\end{mylist}
%	
%	
%%	\item
%%	Set  \( \PEDR_{\jl}^{(0)} \) as solution \( \PEDR_{\jl} \) from the last step.
%	
%	\item
%	Define the set \( \JJJd = \{ \jl \} \).
%
%\end{mylist}
%
%\item
%Take the next point \( \xv_{\jl} \) with \( \jl \not\in \JJJd \).
%
%\begin{mylist}
%	\item
%	Compute \( \BEDR_{\jl}^{(0)} \) by extrapolation
%\begin{EQA}[c]
%	\BEDR_{\jl}^{(0)}
%	\eqdef
%	\biggl( \sum_{j \in \JJJd} \weightM_{\jjl} \biggr)^{-1} \sum_{j \in \JJJd} \PEDR_{j}^{(0)} \; \weightM_{\jjl}
%%	\, ,
%%	\qquad
%%	\WeightM_{\jl}
%%	\eqdef
%%	\sum_{j \in \JJJd} \weightM_{\jjl}	
%	\, .
%\label{jfu88uuujecvyyyyyued}
%\end{EQA}
%
%	\item 
%	Fix \( \BEDR_{\jl} \) by solving \eqref{jjhidkmnvuy5rrtvnv}.
%%\begin{EQA}[c]
%%	\BEDR_{\jl}
%%	=
%%	\arginf_{\BEDR \in \R^{\dimm \times \dimd}} \biggl( \sum_{j \in \JJJ}
%%		\bigl\| \Imav_{j} - \UV_{j} \BEDR^{\T} \fclv_{\jjl} \bigr\|^{2} \weightM_{\jjl}
%%		+ \lambda_{\subM} \| \BEDR  - \BEDR_{\jl}^{(0)} \|_{\Fr}^{2} 
%%	\biggr)
%%	\, .
%%	\qquad
%%\end{EQA}
%
%	\item
%	Compute \( \PEDR_{\jl} \) and \( \Lambda_{\jl} \) by SVD of \( \JEDR_{\jl} \) as in \eqref{jhuifdyu8v8i6445hve} and
%	\eqref{jhfuui887678ucfhie3}.
%
%	\item 
%	Set \( \PEDR_{\jl}^{(0)} = \PEDR_{\jl} \) and \( \Lambda_{\jl}^{(0)} = \Lambda_{\jl} \).
%	
%	\item
%	Add \( \jl \) to \( \JJJd \) and loop until \( \JJJd = \JJJ \).
%%The procedure also delivers \( \{ \JEDR_{\jl} \}_{\jl \in \JJJ} \).
%
%\end{mylist}
%\end{mylist}


Now we describe the step \( k \) of the main procedure for \( k \geq 0 \).


\paragraph{Step \( k \). Tensors
\( (\TEDR_{j}^{(k)} )_{j \in \JJJ} \) and preliminary guess \( (\PEDR_{j,\init}^{(k)} )_{j \in \JJJ} \) are available.}

\begin{mylist}


\item
For each \( j \in \JJJ \),
generate sets \( \Sph_{j} \) of random directions  
\( \sph \eqd \mfrac{\gaussv}{\| \gaussv \|} \), 
\( \gaussv \sim \ND(0,\Id_{\dimd}) \);

\item 
For each \( j \in \JJJ \), compute 
\( \neff_{j}^{(k)} \), \( \bigl( \Ima_{j,\sph}^{(k)} \), \( \Uv_{j,\sph}^{(k)} \bigr)_{\sph \in \Sph_{j}} \)
by \eqref{uyhvvy883hierwfewAVE} and \eqref{juu8hjfioee433erhAVE} with
\begin{EQA}[c]
	\weight_{ij} 
	= 
	\weight_{ij}^{(k)}
	= 
	\Kern\bigl( \| \TEDR_{k} (\Xv_{i} - \xv_{j}) \|^{2} \bigr) 
	\, .
\end{EQA}
Define 
\( \Imav_{j}^{(k)} = \bigl( \Ima_{j,\sph}^{(k)} \bigr)_{\sph \in \Sph_{j}} \),
%\( \Sv_{j}^{(k)} = \bigl( S_{j,\sph}^{(k)} \bigr)_{\sph \in \Sph_{j}} \),
\( \UV_{j}^{(k)} = {(\Uv_{j,\sph}^{(k)})}_{\sph \in \Sph_{j}}^{\T} \).

\item
For each \( \jl \in \JJJ \):
\begin{mylist}
	\item
	Fix \( (\fclv_{\jjl}^{(k)})_{j \in \JJJ} \) by solving 
\begin{EQA}[c]
	\fclv_{\jjl}^{(k)}
	\eqdef
	\arginf_{\fclv_{\jjl}}
	\bigl\| \Imav_{j}^{(k)} - \UV_{j}^{(k)} {\PEDR_{\jl}^{(k-1)}}^{\T} \fclv_{\jjl} \bigr\|^{2}
	\, ,
	\quad
	j \in \JJJ
	\, .
	\qquad
\label{kiiilfduy543rbgckr7r}
\end{EQA}
%	\eqref{kiiilfduy543rbggioe46}
%	with \( \Imav_{j}^{(k)} \, , \Sv_{j}^{(k)} \, , \UV_{j}^{(k)} \) in place of 
%	\( \Imav_{j} \, , \Sv_{j} \, , \UV_{j} \).
%
	\item
	Compute \( \IEDR_{\jl}^{(k-1)} \) as
\begin{EQA}[c]
	\IEDR_{\jl}^{(k-1)}
	\eqdef
	\biggl( \sum_{j \in \JJJ} \weightM_{\jjl}^{(k)} \biggr)^{-1} 
	\sum_{j \in \JJJ} \bigl( \PEDR_{j}^{(k-1)} \bigr)^{\T} \PEDR_{j}^{(k-1)} \; \weightM_{\jjl}^{(k)}
%	\, ,
%	\qquad
%	\WeightM_{\jl}
%	\eqdef
%	\sum_{j \in \JJJ} \weightM_{\jjl}	
	\, .
\label{jfu88uuujpvue2rchuu}
\end{EQA}

%
	\item
	Solve
\begin{EQA}[c]
	\BEDR_{\jl}^{(k)}
	=
	\arginf_{\BEDR \in \R^{\dimm \times \dimd}} \biggl( \sum_{j \in \JJJ}
		\bigl\| \Imav_{j}^{(k)} - \UV_{j}^{(k)} \BEDR^{\T} \fclv_{\jjl}^{(k)} \bigr\|^{2} \neff_{j}^{(k)} \, 
		\weightM_{\jjl}^{(k)}
		+ \lambda_{\subM} \tr\bigl\{ \BEDR^{\T} \BEDR (\Id_{\dimd}  - \IEDR_{\jl}^{(k-1)} ) \bigr\}
	\biggr)
	\, .
	\qquad
\end{EQA}
%
	\item
	Compute \( \JEDR_{\jl}^{(k)} \) by \eqref{jhuifdyu8v8i6445hve} and 
	\( \PEDR_{\jl}^{(k)} \, , \Lambda_{\jl}^{(k)} \) by \eqref{jhfuui887678ucfhie3}
	with \( \BEDR_{\jl} = \BEDR_{\jl}^{(k)} \), \( \fclv_{\jjl} = \fclv_{\jjl}^{(k)} \).

\end{mylist}

%\item 
%Compute the data fit
%\begin{EQA}[c]
%	\err_{k}
%	\eqdef
%	\sum_{j \in \JJJ} \sumi 
%	\bigl\{ Y_{i} - S_{j}^{(k)} - \fcl_{j}^{(k)} (\Xv_{i} - \Xbar_{j}^{(k)})^{\T} \betav_{k} \bigr\}^{2} \, 
%	\weight_{ij}^{(k)}
%	\, ,
%\end{EQA}
%where
%\begin{EQA}[c]
%	\Xbar_{j}^{(k)}
%	=
%	\biggl( \sumi \weight_{ij}^{(k)} \biggr)^{-1} \sumi \Xv_{i} \, \weight_{ij}^{(k)}
%	\, ,
%	\qquad
%	S_{j}^{(k)}
%	=
%	\biggl( \sumi \weight_{ij}^{(k)} \biggr)^{-1} \sumi Y_{i} \, \weight_{ij}^{(k)}
%	\, .
%\end{EQA}

\item
If \( \hiso_{k}/a < \hiso_{\min} \) then terminate with output \( \PEDR_{j} = \PEDR_{j}^{(k)} \).

\item
Otherwise

\begin{mylist}
\item
Update \( \PEDR_{j,\init} = \PEDR_{j}^{(k)} \), \( j \in \JJJ \).

\item 
Increase \( k \) by one.
Update \( \hiso_{k} = \hiso_{k-1} / a \).

\item
With \( \JEDR_{\jl}^{(k-1)} = {\PEDR_{\jl}^{(k-1)}}^{\T} \Lambda_{\jl}^{(k-1)} \, \PEDR_{\jl}^{(k-1)} \),
define \( \TEDR_{j}^{(k)} \) and \( \TEDRm_{\jl}^{(k)} \) by
\begin{EQA}[rcl]
	(\TEDR_{j}^{(k)})^{2}
	&=&
	\frac{\aniso_{k}^{2}}{\hiso_{k}^{2}} \bigl( \Id_{\dimd} - {\PEDR_{j}^{(k-1)}}^{\T} \PEDR_{j}^{(k-1)} \bigr) 
	+ \frac{1}{\hiso_{k}^{2}} \, \JEDR_{j}^{(k-1)}
	\, ,
	\qquad
	j \in \JJJ
	\, ,
	\\
	(\TEDRm_{\jl}^{(k)})^{2}
	&=&
	\frac{\anisom_{k}^{2}}{\hisom_{k}^{2}} \bigl( \Id_{\dimd} - {\PEDR_{\jl}^{(k-1)}}^{\T} \PEDR_{\jl}^{(k-1)} \bigr) 
	+ \frac{1}{\hisom_{k}^{2}} \, \JEDR_{\jl}^{(k-1)}
	\, ,
	\qquad
	\jl \in \JJJ
	\, ,
\end{EQA}
where the anisotropy indices \( \aniso_{k} \) and \( \anisom_{k} \) are the largest values ensuring
\begin{EQA}[rcl]
	\frac{1}{\nJJJ} \sum_{j \in \JJJ} \sumi \Kern\bigl( \| \TEDR_{j}^{(k)} (\Xv_{i} - \xv_{j}) \|^{2} \bigr)
	& \geq &
	\nmin
	\, ,
	\\
	\frac{1}{\nJJJ} \sum_{\jl \in \JJJ} \sum_{j \in \JJJ} \Kern\bigl( \| \TEDRm_{\jl}^{(k)} (\xv_{j} - \xv_{\jl}) \|^{2} \bigr)
	& \geq &
	\nman
	\, .
\label{jhgiverkytde4w34etycjamM}
\end{EQA}
\end{mylist}

%
	\item
	Loop.

\end{mylist}

{\color{blue} please, evaluate
\begin{mylist}

	\item 
	memory requirement as function of \( n \), \( \dimd \), \( \dimm \), \( \nJJJ \), \( \nSph \);
	\item
	computational cost as function of \( n \), \( \dimd \), \( \dimm \), \( \nJJJ \), \( \nSph \);
	\item
	accuracy of estimating the gradients \( \nabla \gs(\xv_{\jl}) \);
	\item
	accuracy of estimation after initialization (after the first run and then after each repetition)
	and frequency of a failure;
	\item
	one-step improvement by structural adaptation;
	\item
	choice of \( \lambda_{\manifold} \);
	\item
	degradation with growth of \( \dimd \).
\end{mylist}

}
